\chapter*{How to Read This Collection}
\addcontentsline{toc}{chapter}{How to Read This Collection}

\section*{The plain-terms creation story}

Each tradition is introduced by its creation story told in ordinary modern words, with no native terms at all.
Instead of naming a god, we name the role that god plays, the source, the first stirring, the dividing of the
realms, the breaking, the spark, the return. This stripped retelling is the main device of the collection. It
lets a reader who knows none of the traditions see the structure at once, and it lets a reader who knows one
tradition see how its familiar story is also being told, in different form, everywhere else.

The same creation story appears twice. Here in the body, in plain terms. And again in the appendix, retold in
the tradition's own vocabulary, so the two can be compared directly.

\section*{The mapping of inner-experience language}

Underneath the surface differences, the traditions are often pointing at the same inner realities, the ground
of being, the witness behind thought, the breath that is also spirit, the light that is also knowledge, the
descent and the return. This collection uses a small, consistent set of plain modern phrases for these
realities, and then maps each tradition's native term onto them.

For example, the plain phrase \textit{the unknowable source} maps to \textbf{Ein Sof} in Kabbalah, to
\textbf{emptiness} or \textbf{the Ground} in Buddhism, to \textbf{Brahman} in Hinduism, to the \textbf{Tao}
in Taoism, to \textbf{the One} in Neoplatonism, and to the \textbf{Essence} in Sufism. The plain phrase
\textit{the breath that is spirit} maps to \textbf{ruach}, \textbf{prana}, \textbf{qi}, \textbf{pneuma}, and
\textbf{nafas}. The full mappings, including native script such as Hebrew where it adds value and a standard
romanization first, are gathered in the appendix.

\section*{On tables and on prose}

Many traditions organize reality into numbered sets, the ten emanations, the five elements, the seven heavens,
the eighteen conditions of a precious human life. Where a tradition counts, this collection uses a table,
because a table shows a structure better than a paragraph. But the collection is not only tables. The prose
carries the story, the meaning, and the comparison, and the tables sit inside it as the underlying structure.

\section*{On comparison}

Each tradition is compared, as it is introduced, with the ones before it. Each major part ends with a section
that draws its traditions together. And the whole collection ends with a synthesis that proposes the single
model the many are circling. The comparisons are the point. Any one tradition, read alone, looks like a closed
world. Read against the others, it becomes one variant of a shared structure that recurs across cultures.

\section*{On accuracy and honesty}

The traditions disagree internally, and their texts are sometimes fragmentary or filtered through later hands.
Where a claim is contested, a reconstruction, or version-dependent, the text says so plainly rather than
smoothing it over. The goal is a faithful map, including its uncertain regions, not a tidy fiction.
